\SourceLeaf{275}
\[
\begin{aligned}
Q+q&=46^0\ 40'\ 37''69794\\
Q+q&=58\quad39\quad44{,}09285
\end{aligned}
\]
while the finite formula gives
\[
\begin{aligned}
Q+q&=46^0\ 40'\ 37''69794\\
Q+q&=58\quad39\quad44{,}09283
\end{aligned}
\]
so the results agree as closely as ten-digit logarithms allow.

\Article{7}
\indent Similarly, the logarithm of \(m\) can be expanded in a series whose first terms are as follows.
\[
\begin{aligned}
\log m={}&-\frac{\sin2\varphi^2}{6\cos\theta^4}\mathord{.}csp^3-\frac{\sin2\varphi^2}{24\cos\theta^6}(cc+11eess)p^4\\
&+\frac{\sin2\varphi^2}{120\cos\theta^8}\mathord{.}\frac{s}{c}(2cc-3ss-ee(40c^4-20ccss-6s^4)\\
&\qquad-e^4ss(104c^4+22ccss+3s^4))p^5
\end{aligned}
\]
I have also worked out the next term (by a different method), though only as far as the principal part of its coefficient, which is of order \(ee\). I obtained
\[
+\frac{\sin2\varphi^2}{720\cos\theta^{10}}\mathord{.}\frac{1}{cc}(2c^4-18ccss-15s^4)p^6
\]
\indent The logarithm expressed by this series is the hyperbolic logarithm, and \(p\) is understood, as above, to be expressed in parts of the radius. If the Briggsian logarithm is wanted, with \(p\) denoting degrees, the modulus must be included as a factor and \(\tfrac{\pi p}{180}\) written in place of \(p\). In this form our example gives
\[
\begin{aligned}
\log m={}&-0{,}0049612433(\frac{p}{100})^3\\
&-0{,}0017329876(\frac{p}{100})^4\\
&-0{,}002393772(\frac{p}{100})^5\\
&-0{,}0124746(\frac{p}{100})^6
\end{aligned}
\]
